%% Le lingue utilizzate, che verranno passate come opzioni al pacchetto babel. Come sempre, l'ultima indicata sarà quella primaria.
%% Se si utilizzano una o più lingue diverse da "italian" o "english", leggere le istruzioni in fondo.
\def\thudbabelopt{english,italian}
%% Valori ammessi per target: bach (tesi triennale), mst (tesi magistrale), phd (tesi di dottorato).
%% Valori ammessi per aauheader: '' (vuoto -> nessun header Alpen Adria Univeristat), aics (Department of Artificial Intelligence and Cybersecurity), informatics (Department of Informatics Systems). Il nome del dipartimento è allineato con la versione inglese del logo UniUD.
%% Valori ammessi per style: '' (vuoto -> stile moderno), old (stile tradizionale).
\documentclass[target=bach,aauheader=,style=]{thud}

%% --- Informazioni sulla tesi ---
%% Per tutti i tipi di tesi
% Scommentare quello di interesse, o mettete quello che vi pare
\course{Informatica}
%\course{Internet of Things, Big Data e Web}
%\course{Matematica}
%\course{Comunicazione Multimediale e Tecnologie dell'Informazione}
\title{Utilizzo di telecamere \\ in ambienti industriali reali}
\author{Michele Ongaro}
\supervisor{Prof.\ Agostino Dovier}
% \cosupervisor{Arch.\ Rambaldo Melandri \and Dott.\ Giorgio Perozzi}
\tutor{Ing. Umberto Piconi}
%% Campi obbligatori: \title, \author e \course.
%% Altri campi disponibili: \reviewer, \tutor, \chair, \date (anno accademico, calcolato in automatico), \rights
%% Con \supervisor, \cosupervisor, \reviewer e \tutor si possono indicare più nomi separati da \and.
%% Per le sole tesi di dottorato:
% \phdnumber{313}
% \cycle{XXVIII}
% \contacts{Via della Sintassi Astratta, 0/1\\65536 Gigatera --- Italia\\+39 0123 456789\\\texttt{http://www.example.com}\\\texttt{inbox@example.com}}

%% --- Pacchetti consigliati ---
%% pdfx: per generare il PDF/A per l'archiviazione. Necessario solo per la versione finale
\usepackage[a-1b]{pdfx}
%% hyperref: Regola le impostazioni della creazione del PDF... più tante altre cose. Ricordarsi di usare l'opzione pdfa.
\usepackage[pdfa]{hyperref}
%% tocbibind: Inserisce nell'indice anche la lista delle figure, la bibliografia, ecc.

%% --- Stili di pagina disponibili (comando \pagestyle) ---
%% sfbig (predefinito): Apertura delle parti e dei capitoli col numero grande; titoli delle parti e dei capitoli e intestazioni di pagina in sans serif.
%% big: Come "sfbig", solo serif.
%% plain: Apertura delle parti e dei capitoli tradizionali di LaTeX; intestazioni di pagina come "big".

\begin{document}
\maketitle

%% Dedica (opzionale)
% \begin{dedication}
% 	Al mio cane,\par per avermi ascoltato mentre ripassavo le lezioni.
% \end{dedication}

%% Ringraziamenti (opzionali)
% \acknowledgements
% Sed vel lorem a arcu faucibus aliquet eu semper tortor. Aliquam dolor lacus, semper vitae ligula sed, blandit iaculis leo. Nam pharetra lobortis leo nec auctor. Pellentesque habitant morbi tristique senectus et netus et malesuada fames ac turpis egestas. Fusce ac risus pulvinar, congue eros non, interdum metus. Mauris tincidunt neque et aliquam imperdiet. Aenean ac tellus id nibh pellentesque pulvinar ut eu lacus. Proin tempor facilisis tortor, et hendrerit purus commodo laoreet. Quisque sed augue id ligula consectetur adipiscing. Vestibulum libero metus, lacinia ac vestibulum eu, varius non arcu. Nam et gravida velit.

%% Sommario (opzionale)
% \abstract
% Nunc ac dignissim ipsum, quis pulvinar elit. Mauris congue nec leo ornare lobortis. Nulla hendrerit pretium diam nec lobortis. Nullam aliquam laoreet nisl, sit amet facilisis lectus accumsan ut. Duis et elit hendrerit metus venenatis condimentum. Integer id eros molestie, interdum leo sit amet, aliquet metus. Integer fermentum tristique magna, vel luctus neque rhoncus vel. Ut hendrerit et quam et semper. Mauris egestas, odio sed aliquet luctus, magna orci euismod odio, vitae lacinia tellus tellus non lectus. Aliquam urna neque, porta et mattis aliquam, congue sit amet lorem. In ultrices augue sit amet ante vehicula, vitae rhoncus turpis auctor. Donec porta scelerisque eros, at mollis enim imperdiet ut. 

%% Indice
\tableofcontents

%% Lista delle tabelle (se presenti)
%\listoftables

%% Lista delle figure (se presenti)
%\listoffigures

%% Corpo principale del documento
\mainmatter

%% Parte
%% La suddivisione in parti è opzionale; solitamente sono sufficienti i capitoli.
%\part{Parte}

%% Capitolo
\chapter{Introduzione}

\section{Obiettivi}

Si vuole realizzare un algoritmo che permetta di determinare se l'obiettivo di una telecamera sia sporco o meno tramite l'analisi delle immagini acquisite. In particolare, le immagini da considerare provengono da un flusso video ottenuto tramite delle telecamere IP posizionate in corrispondenza di una serie di targhe (\textit{plate}) contenenti un codice realizzato tramite una codifica Hamming.

\section{Presentazione scenario applicativo}

L'infrastruttura aziendale presa in considerazione prevede una serie di telecamera IP, posizionate in punti strategici, raffiguranti delle cisterne contenenti materiale metallico fuso e con su appesa una targa (\textit{plate}) contenente un certo codice. Tramite protocollo RTSP, le telecamere, forniscono un flusso di immagini ad una \textit{pipeline} di riconoscimento di codici. Tali immagini vengono processate dalla \textit{pipeline}, che attraverso una serie livelli di elaborazione riesce ad ottenere il valore numerico associato al codice del \textit{plate}. Tuttavia, non sempre per la \textit{pipeline} è possibile riconoscere tale valore, e questo è dovuto a due motivi principali. Il primo è dato dal fatto che il \textit{plate} non è affatto presente nell'immagine o nell'area presa in considerazione dalla \textit{pipeline}, infatti la procedura che porta al riconoscimento del valore numerico associato al \textit{plate} può essere configurata in modo tale da considerare esclusivamente una sotto-area dell'immagine originale. Il secondo motivo per cui la pipeline non è capace di riconoscere il valore numerico è legato a fenomeni esterni come una scarsa illuminazione, una messa a fuoco dell'immagine precaria, la presenza parziale del \textit{plate} o la stessa dimensione che il quest'ultimo assume nell'immagine complessiva. Un'altra casistica che porta la \textit{pipeline} a fallire è la presenza di sporcizia sull'obiettivo, capita infatti che delle cisterne durante il loro movimento degli schizzi di metallo cadano sull'obiettivo della telecamera andandolo a sporcare. Il riconoscitore viene configurato con un certo limite di affidabilità e in questi casi succede che tale livello non può essere garantito dalla decodifica del codice e ciò comporta un esito negativo da parte dell'intera \textit{pipeline}.

Per questo motivo è importante poter far affidamento ad un indice che sia indipendente dalla presenza del \textit{plate} all'interno dell'immagine che permetta di valutare se l'obiettivo della telecamera sia in condizioni ottimali. In questo modo sarebbe possibile correggere la ripresa di una telecamera in anticipo e prevenire lunghi periodi di tempo durante i quali la \textit{pipeline} di riconoscimento non valuta alcun valore, ma non perché il codice è assente nell'immagine ma perché l'obiettivo della telecamera non ha la possibilità di fare delle riprese pulite.

\section{Considerazioni implementative}

Le telecamere sono configurate in modo tale da gestire la messa a fuoco in maniera automatica. Si è osservato che nei casi in cui l'obiettivo della telecamera viene sporcato da schizzi di metallo, la messa a fuoco cambia e prende come riferimento la macchia presente sull'obiettivo, questo ha come effetto che la restante porzione dell'immagine diventi completamente sfocata e la \textit{pipeline} di riconoscimento non garantisce più un adeguato livello di affidabilità. Per questo motivo lo studio che verrà illustrato di seguito si basa esclusivamente sulla considerazione di immagini sfocate / nitide. Se quindi l'algoritmo riuscirà a trovare che una immagine è sfocata allora questo rappresentarà un indice chiaro del fatto che la telecamera è fuori fuoco e la causa di ciò sarà con buona probabilità semplicemente dovuto al fatto che l'obiettivo è stato sporcato.

\chapter{Sviluppo delle metriche}

Per ottenere un indice che permetta di determinare la nitidezza di una immagine si può procedere facendo uso di una serie di approcci diversi. Si può fare distinzione tra le \textbf{metriche basate sull'analisi delle frequenze} e le \textbf{metriche basate sull'\textit{edge detection}}. Queste sono le due tipologie principali ma non sono le sole, infatti si può andare ad osservare ulteriori caratteristiche della natura dell'immagine che non sono, almeno direttamente, associabili ad una di queste due macro tipologie.


\section{Metriche basate sull'analisi delle frequenze}

L'idea che sta alla base dell'utilizzo di una metrica basata sull'analisi delle frequenze, sta nel considerare una immagine come un segnale bidimensionale. Le immagini infatti sono dei segnali discreti a due dimensioni. Il valore assunto da ciascun \textit{pixel} è paragonabile al valore assunto dall'ordinata in un segnale monodimensionale. Se quindi, utilizzando strumenti come la trasformata di Fourier, è possibile estrarre lo spettro delle frequenze, allora si può constatare che una immagine nitida, rispetto ad una sfocata, sarà composta da frequenze più elevate. Tali frequenze nell'immagine vengono rappresentate dai piccoli dettagli o dalle rapide variazioni di colore, caratteristica invece non presente in una immagine sfocata.

Entrambe le metriche che seguono si basano sulla trasformata discreta di Fourier (DTF) e differiscono nel modo in cui tale risultato viene valutato.

\subsection{Fourier}

\subsection{Skewness}


\section{Metriche basate sull'edge detection}

Per quanto riguarda le metriche basate sull'\textit{edge detection} il concetto alla base delle diverse interpretazioni è che un'immagine nitida ha dei bordi ben definiti e in quantità maggiore rispetto ad un'immagine sfocata.

\subsection{Laplacian}

\subsection{Canny}

\section{Altre metriche}

\chapter{Definizione dei threshold}

Qui spiego come sono stati definiti i boundary per ogni metrica.

\chapter{Valutazione e confronto delle metriche}



\chapter{Conclusioni e considerazioni finali}


%% Fine dei capitoli normali, inizio dei capitoli-appendice (opzionali)
\appendix

%\part{Appendici}

\chapter{Titolo della prima appendice}
Sed purus libero, vestibulum ut nibh vitae, mollis ultricies augue. Pellentesque velit libero, tempor sed pulvinar non, fermentum eu leo. Duis posuere eleifend nulla eget sagittis. Nam laoreet accumsan rutrum. Interdum et malesuada fames ac ante ipsum primis in faucibus. Curabitur eget libero quis leo porttitor vehicula eget nec odio. Proin euismod interdum ligula non ultricies. Maecenas sit amet accumsan sapien.

%% Parte conclusiva del documento; tipicamente per riassunto, bibliografia e/o indice analitico.
\backmatter

%% Riassunto (opzionale)
%\summary
%Maecenas tempor elit sed arcu commodo, dapibus sagittis leo egestas. Praesent at ultrices urna. Integer et nibh in augue mollis facilisis sit amet eget magna. Fusce at porttitor sapien. Phasellus imperdiet, felis et molestie vulputate, mauris sapien tincidunt justo, in lacinia velit nisi nec ipsum. Duis elementum pharetra lorem, ut pellentesque nulla congue et. Sed eu venenatis tellus, pharetra cursus felis. Sed et luctus nunc. Aenean commodo, neque a aliquam bibendum, mauris augue fringilla justo, et scelerisque odio mi sit amet diam. Nulla at placerat nibh, nec rutrum urna. Donec ut egestas magna. Aliquam erat volutpat. Phasellus vestibulum justo sed purus mattis, vitae lacinia magna viverra. Nulla rutrum diam dui, vel semper mi mattis ac. Vestibulum ante ipsum primis in faucibus orci luctus et ultrices posuere cubilia Curae; Donec id vestibulum lectus, eget tristique est.

%% Bibliografia (praticamente obbligatoria)
\bibliographystyle{plain_\languagename}%% Carica l'omonimo file .bst, dove \languagename è la lingua attiva.
%% Nel caso in cui si usi un file .bib (consigliato)
\bibliography{thud}
%% Nel caso di bibliografia manuale, usare l'environment thebibliography.

%% Per l'indice analitico, usare il pacchetto makeidx (o analogo).

\end{document}

--- Istruzioni per l'aggiunta di nuove lingue ---
Per ogni nuova lingua utilizzata aggiungere nel preambolo il seguente spezzone:
    \addto\captionsitalian{%
        \def\abstractname{Sommario}%
        \def\acknowledgementsname{Ringraziamenti}%
        \def\authorcontactsname{Contatti dell'autore}%
        \def\candidatename{Candidato}%
        \def\chairname{Direttore}%
        \def\conclusionsname{Conclusioni}%
        \def\cosupervisorname{Co-relatore}%
        \def\cosupervisorsname{Co-relatori}%
        \def\cyclename{Ciclo}%
        \def\datename{Anno accademico}%
        \def\indexname{Indice analitico}%
        \def\institutecontactsname{Contatti dell'Istituto}%
        \def\introductionname{Introduzione}%
        \def\prefacename{Prefazione}%
        \def\reviewername{Controrelatore}%
        \def\reviewersname{Controrelatori}%
        %% Anno accademico
        \def\shortdatename{A.A.}%
        \def\summaryname{Riassunto}%
        \def\supervisorname{Relatore}%
        \def\supervisorsname{Relatori}%
        \def\thesisname{Tesi di \expandafter\ifcase\csname thud@target\endcsname Laurea\or Laurea Magistrale\or Dottorato\fi}%
        \def\tutorname{Tutor aziendale}%
        \def\tutorsname{Tutor aziendali}%
    }
sostituendo a "italian" (nella 1a riga) il nome della lingua e traducendo le varie voci.
